Mit seiner sogenannten «Mondrechnung» zeigte Newton 1666, dass die Gravitationskraft zwischen Erde und Mond die gleiche Kraft ist, die den freien Fall eines Apfels auf der Erde verursacht.

\begin{abcliste}
\abc Berechnen Sie die Radialbeschleunigung des Mondes aus seiner mittleren Entfernung $r_\mathrm{EM}$ und seiner Umlaufzeit um die Erde.
\abc Zeigen Sie, dass für das Verhältnis der Fallbeschleunigung auf der Erde (Erdradius $r_\mathrm{E}$) zur Radialbeschleunigung des Mondes gilt:
\begin{align}
 g : a_R &= \left(r_\mathrm{EM} : r_\mathrm{E}\right)^2
\end{align}
\abc Vergleichen Sie die Resultate von a) und b) und kommentieren Sie Ihre Überlegungen.
\end{abcliste}
